% Section 4.3, from lectures/L04/README.md: the objectives, the evaluation questions, and the next
% lecture.

\needspace{10\baselineskip}
\section{Review}
\label{c4:sec:review}

\subsection*{What you should be able to do now}
After this chapter, you should:
\begin{itemize}
  \item Have created an interface for a neural network.
  \item Be able to connect components via interfaces.
  \item Have implemented a concrete class for a simple neural network, including a fully working
    training method.
  \item Be able to explain how a model can set its own learning rate, and why a training method
    that resets its parameters first is worth having.
\end{itemize}

\needspace{6\baselineskip}
\subsection*{Questions to test yourself}
You should be able to explain:
\begin{enumerate}
  \item Why is it beneficial to define an interface for the neural network rather than
    implementing a class directly?
  \item Without looking at the code, can you describe how the components in the network are
    connected via interfaces?
  \item Can you explain the flow from input to prediction in your implementation, step by step?
  \item Why doesn't \code{feedforward()} return the layer's output, and how does the rest of the
    network access it instead?
  \item Why is the precision evaluated every hundredth epoch here, rather than every tenth as in
    Chapter~\ref{c2:ch:training}, and what does the network do with the number?
  \item Why does \code{train()} reset both layers before its first epoch, and what would calling it
    twice do without that reset?
\end{enumerate}

\needspace{6\baselineskip}
\subsection*{Looking ahead}
Chapter~\ref{c5:ch:dense} implements a real dense layer in software, replacing the stub used since
Chapter~\ref{c3:ch:nn}.
